\documentclass[12pt]{article}
\usepackage[T1]{fontenc}
\usepackage[utf8]{inputenc}
\usepackage{lmodern}
\usepackage{textcomp}
\usepackage{enumitem}
\usepackage{geometry}
\geometry{margin=1in}
\begin{document}
\begin{center}
Answer Key - Political Science - Graduate Level Assessment 23 \
\small (Answers are ordered exactly as in the question paper.)
\end{center}

\section*{Multiple Choice}
\begin{enumerate}[label=\arabic*.]
\item \textbf{Answer:} D
\\ \textit{Options:} A) The nexus between identity and security.; B) War, the military, and the sovereign state.; C) Environmental changes.; D) All of these options.
\\ \textit{Note:} The correct answer is "All of these options" because the 'deepening and broadening' of security encompasses various dimensions, including traditional military concerns, economic stability, environmental sustainability, and social cohesion, all of which contribute to a more comprehensive understanding of security in contemporary political science.
\item \textbf{Answer:} A
\\ \textit{Options:} A) Business networks; B) Computer experts; C) Anti-virus industry; D) Computer networks
\\ \textit{Note:} Business networks do not form part of the technical discourse on cybersecurity because they primarily focus on organizational structures and relationships rather than the technical measures and protocols used to protect information systems.
\item \textbf{Answer:} C
\\ \textit{Options:} A) Ronald Reagan; B) George Soros; C) Bill Clinton; D) George W. Bush
\\ \textit{Note:} Bill Clinton articulated the idea that globalization is an inevitable and powerful force in the modern economy, likening it to natural elements, which reflects his perspective on the interconnectedness of global markets during his presidency.
\item \textbf{Answer:} B
\\ \textit{Options:} A) Digitalised sensitive information; B) Critical Information Infrastructures; C) Government IT systems; D) Telecommunication networks
\\ \textit{Note:} The correct answer is Critical Information Infrastructures because they are essential systems and assets whose disruption or compromise would have a significant impact on national security, economic stability, and public safety, making them the primary focus of contemporary cyber-security efforts.
\item \textbf{Answer:} A
\\ \textit{Options:} A) Provide strong enforcement; B) Monitor parties; C) Provide fora for discussion; D) Reduce transaction costs for agreements
\\ \textit{Note:} Intergovernmental organizations seldom provide strong enforcement because they primarily rely on member states to implement agreements and lack direct authority or coercive power to impose compliance.
\end{enumerate}

\section*{True / False}
\begin{enumerate}[label=\arabic*.]
\setcounter{enumi}{5}
\item \textbf{Answer:} True
\\ \textit{Options:} A) True; B) False
\\ \textit{Note:} The statement is true because the foreign policies of other states, international law, and intergovernmental organizations significantly influence and limit the options available to US policymakers, shaping the context within which foreign policy decisions are made.
\item \textbf{Answer:} False
\\ \textit{Options:} A) True; B) False
\\ \textit{Note:} Rosenau argues that U.S. foreign policy behavior is influenced by a complex interplay of various factors, including domestic politics, international systems, and historical contexts, rather than solely the individual personality of the President.
\item \textbf{Answer:} False
\\ \textit{Options:} A) True; B) False
\\ \textit{Note:} The Roosevelt Doctrine, articulated by President Theodore Roosevelt, actually asserted the United States' right to intervene in Latin America to maintain stability, thereby diminishing the role of European intervention in the region.
\item \textbf{Answer:} True
\\ \textit{Options:} A) True; B) False
\\ \textit{Note:} Exceptionalism posits that unique national characteristics can justify both isolationist and internationalist approaches, as both aim to promote liberal values such as democracy and human rights, albeit through different strategies.
\item \textbf{Answer:} True
\\ \textit{Options:} A) True; B) False
\\ \textit{Note:} This statement is true because free trade typically leads to competitive advantages for certain industries, resulting in job losses and economic disruptions for those specific sectors, while consumers across the economy benefit from lower prices and increased product variety.
\end{enumerate}

\section*{Long Form Response}
\begin{enumerate}[label=\arabic*.]
\setcounter{enumi}{10}
\item \textbf{Answer:} The human security concept has faced several criticisms that impact its effectiveness in addressing contemporary global issues. Critics argue that the concept is too broad and lacks a clear definition, leading to ambiguous policy applications that can dilute its significance. Additionally, some contend that it shifts focus away from traditional state security, potentially undermining the responsibility of states to protect their citizens. Furthermore, the emphasis on individual security can overlook structural and systemic issues, such as economic inequality and political instability, which require collective action. These criticisms suggest that while human security has the potential to address pressing global challenges, its practical implementation may be hindered by conceptual vagueness and a lack of coherence in prioritizing state versus individual needs.
\\ \textit{Note:} The answer correctly identifies key criticisms of the human security concept, such as its broadness, ambiguity in policy application, potential neglect of state responsibilities, and the overemphasis on individual security, all of which hinder its effectiveness in addressing contemporary global issues.
\item \textbf{Answer:} The interplay between culture and social constructivism is vital in shaping state identity and security perceptions in political science. Culture, understood as a set of practices that imbue shared experiences with meaning, influences how states interpret security threats and define their national interests. Social constructivism posits that identities and interests are not predetermined but are constructed through social interactions, wherein cultural narratives and historical contexts play crucial roles. Thus, the values and norms derived from a state's cultural heritage significantly inform its security policies and identity formation, leading to diverse approaches to security across different states. This interplay illustrates that security perceptions are not merely strategic calculations but are deeply rooted in the cultural identities that states cultivate.
\\ \textit{Note:} correct because it effectively illustrates how culture, through shared meanings and practices, interacts with social constructivism to shape the fluid identities and security perceptions of states, emphasizing that these constructs are formed through social interactions and historical contexts rather than being fixed or predetermined.
\end{enumerate}

\vfill
\noindent\textit{End of Answer Key}
\end{document}